\section{Architecture}
\label{sec:architecture}

The \Zoo{} \DEX{} runs on \Zoo{} D-Chain --- a white-label resell of
the \Lux{} D-Chain template defined by LP-134 \cite{lp-134} --- which
sits inside the \Lux{} chain topology alongside P-, C-, X-, Q-, Z-,
A-, B-, M-, and F-Chains. The white-label model gives \Zoo{} chain-local
control of asset listings, fee policy, market hours, and validator set
while inheriting the \Lux{} consensus, bridge, privacy, and AI
substrate without re-implementation.

\subsection{Why D-Chain (white-label)}

A native trading venue for tokenised securities needs three properties
simultaneously:

\begin{enumerate}[leftmargin=1.4em]
  \item \textbf{Sub-second finality} so that retail order flow does not
        wait for cross-chain confirmations. \Quasar{} 4.0 triple-cert
        finality on \Lux{} \cite{lp-020} delivers $<1.1$s.
  \item \textbf{Compliance perimeter} so that the chain can enforce
        regulatory constraints (KYC, AML, accredited-investor checks)
        without leaking PII into block contents. The \Liquidity{}
        precompile (\S\ref{sec:liquidity-integration}) provides this.
  \item \textbf{Sovereign listing policy} so that \Zoo{} can list
        \Zoo{}-native securities (per-\LLM{} tokens, conservation
        bonds, fractional \NFT{}s) without coordinating with the broad
        \Lux{} validator set on every market opening.
\end{enumerate}

\Lux{} D-Chain template provides the first two; the white-label resell
adds the third by giving \Zoo{} a sovereign chain instance with its own
governance.

\subsection{Position in the Zoo stack}

\begin{longtable}{p{0.30\textwidth}p{0.62\textwidth}}
\toprule
\textbf{Layer} & \textbf{Role} \\
\midrule
\Zoo{} 4.0 L1 \cite{zoo-4-0} & General-purpose \GPU{}-native sovereign L1; \texttt{\$ZOO} gas; per-\LLM{} chain hub. \\
\Zoo{} D-Chain & Specialised D-Chain instance for tokenised securities. White-label of \Lux{} D-Chain template. \\
\Zoo{} \DEX{} & Application contracts on \Zoo{} D-Chain (\AMM{} V2, V3, V4, order book, order types). \\
\Liquidity{} precompile & Compliance perimeter; KYC/AML; broker-dealer hooks; transfer-agent hooks. \\
\bottomrule
\end{longtable}

\subsection{AMM V2 / V3 / V4}

\paragraph{V2 --- constant product ($x \cdot y = k$).}
Pools are pairs of tokenised assets (security pair, security/stablecoin,
security/\$ZOO). LP shares mint as receipts; default fee 30 bps with
per-pool override governance. Uniswap-shaped surface so existing
analytics, audits, and tooling translate.

\paragraph{V3 --- concentrated liquidity.}
Concentrated-liquidity tick-based pools per the Uniswap v3 design
\cite{uniswap-v3}. Multiple fee tiers ($5/30/100$ bps) coexist for
the same pair. Tick crossings, fee accumulation, and range-position
accounting run as \GPU{} kernels per LP-009 \cite{lp-009} and
LP-132 \cite{lp-132}.

\paragraph{V4 --- precompile to \Lux{} \DEX{}.}
V4 is not a separate \AMM{}. V4 is the precompile interface that lets
\Zoo{} D-Chain contracts route swaps to the \Lux{} \DEX{}
\cite{lux-dex} via B-Chain bridge messages. Use V4 when the deepest
liquidity for a pair lives on \Lux{} (e.g., cross-ecosystem assets,
institutional CLOB flow, cross-chain routing).

\paragraph{Optional local Zoo \DEX{} precompile.}
For \Zoo{}-internal use cases that want CLOB semantics without the
B-Chain hop, \Zoo{} ships an in-chain \DEX{} precompile that
implements the same matching primitives as the \Lux{} \DEX{} but
against \Zoo{} D-Chain order books. The precompile shares its kernel
implementation with \Lux{} \DEX{}; only the venue and gas accounting
differ.

\subsection{Determinism}

The \DEX{} runs deterministically on the \GPU{} per LP-132
\cite{lp-132} and the \GPU{}-Residency Invariant of LP-137
\cite{lp-137}: order matching, fee calculation, and tick crossing
all stay on-device for the duration of a block. The block-level
matching engine reads the order book, runs the matching algorithm
(price-time priority for continuous limit book; uniform-price
discovery for batch auctions), updates pool state, and emits the
trade events all in device memory; only the final cert and event
log leave the device.

\subsection{Settlement}

A trade is final when:

\begin{enumerate}[leftmargin=1.4em]
  \item \Zoo{} D-Chain produces a block containing the matched trades.
  \item \Lux{} \Quasar{} 4.0 emits a triple-cert
        ($\BLS{} + \Corona{} + \MLDSA{}$) over the D-Chain block hash
        \cite{lp-020}.
  \item The trade event is recorded in the chain log; transfer-agent
        hooks fire (\S\ref{sec:liquidity-integration}); the
        beneficial owner record is updated.
\end{enumerate}

End-to-end latency target is below 1.1 seconds from order submission
to triple-cert finality. The latency is dominated by the consensus
round; the matching engine itself completes in microseconds for
typical block contents.

\subsection{Cross-chain plumbing}

\begin{itemize}[leftmargin=1.2em]
  \item \textbf{B-Chain (\textsc{lp-134}).} Cross-chain trade routing,
        bridge messages between \Zoo{} D-Chain and \Lux{} C-Chain or
        external chains.
  \item \textbf{F-Chain (\textsc{lp-013}).} Confidential orderflow:
        ciphertext orders, ciphertext fills, encrypted balances per
        Confidential ERC-20.
  \item \textbf{Z-Chain (\textsc{lp-063}).} \ZKP{} verification for
        accreditation proofs, geographic eligibility, source-of-funds
        attestations.
  \item \textbf{A-Chain (\textsc{lp-134}).} TEE-attested AI mining and
        inference receipts that pay validators in \texttt{\$AI}
        alongside the trading-fee share denominated in
        \texttt{\$ZOO}.
  \item \textbf{M-Chain (\textsc{lp-134}).} Validator-set custody and
        key rotation under \Lux{} threshold ceremony rules (LP-017).
\end{itemize}

\subsection{Why this is the right architecture}

A pure-app \DEX{} on \Zoo{} 4.0 L1 would work for \Zoo{}-native pairs
but would not give the chain a clean compliance boundary against
external regulators. A pure routing layer to \Lux{} \DEX{} would be
fast but would not let \Zoo{} list its own securities sovereignly.
A bespoke chain built from scratch would duplicate two years of
\Lux{} consensus, bridge, and privacy work. The white-label D-Chain
is the right shape: sovereign listing and policy on top of a battle-tested
\Lux{} substrate.
